% TOP_PROBLEM: {"id": 3163, "title": "(m,n)-cycle covers", "category_id": 3, "category": {"id": 3, "name": "graph_theory", "display_name": "Graph Theory", "description": "Problems involving graphs, networks, and their properties.", "slug": "graph-theory", "order_index": 3, "created_at": "2026-07-31T15:26:25.671Z"}, "status": "open", "problem_number": "OPG-139", "set_id": 12}
% FIRST_POSTED: 2026-10-01
% ENTRY_KIND: statement_only
% UnsolvedMath ID 3163; source catalogue code OPG-139
% !TeX program = lualatex
% Definition and sources only; this file is not an LLM solution attempt.
\documentclass[11pt,a4paper]{article}
\usepackage[margin=25mm]{geometry}
\usepackage{fontspec}
\setmainfont{lmroman10-regular.otf}[BoldFont=lmroman10-bold.otf,
 ItalicFont=lmroman10-italic.otf,BoldItalicFont=lmroman10-bolditalic.otf]
\usepackage{amsmath,amssymb,mathtools}
\usepackage{unicode-math}
\setmathfont{latinmodern-math.otf}
\AtBeginDocument{\let\setminus\smallsetminus}
\usepackage{xurl,hyperref}
\hypersetup{colorlinks=true,urlcolor=blue}
\setcounter{secnumdepth}{0}
\setlength{\parindent}{0pt}
\setlength{\parskip}{5pt}
\setlength{\emergencystretch}{3em}
\begin{document}
\section*{(m,n)-cycle covers}\label{3163}
{\small UnsolvedMath ID 3163 · OPG-139 · Graph Theory · OpenGarden}
\subsection{Definitions and mathematical statement}
Let $G=(V,E)$ be a finite undirected loopless multigraph.
A binary cycle is an edge set $C\subseteq E$ such that every
vertex of the spanning subgraph $(V,C)$ has even degree.
The empty set is allowed. A binary cycle may be disconnected
and may have vertices of degree greater than two; equivalently,
its edges can be partitioned into ordinary circuits.

For positive integers $m,n$, an $(m,n)$-cycle cover is a list
$C_1,\ldots,C_m$ of binary cycles such that
\[
 \sum_{i=1}^m\mathbf1_{\{e\in C_i\}}=n
 \qquad\text{for every }e\in E.
\]
The list may contain repeated entries, and empty entries allow
a cover with fewer than $m$ nonempty members to be padded.
Membership of one edge in one binary cycle counts once.

The conjecture is that every bridgeless $G$ admits a
$(5,2)$-cycle cover, where bridgeless means that deleting any
single edge does not increase the number of connected
components. Thus five even-degree edge sets should together
cover each original edge exactly twice.

The number five bounds these binary cycles, not the number
of connected circuits arising after they are decomposed.
In a cubic graph every nonzero degree inside a binary cycle
is two, so each of the five members is a vertex-disjoint union
of circuits. No connectivity of $G$ or of the covering members
is required; components can be handled with the same five
indices.
\subsection{Short English statement}
Can the edges of every bridgeless graph be covered exactly twice by five even-degree subgraphs?
\subsection{Sources}
\begin{itemize}
\item[S1] ULAM.AI and the UnsolvedMath Contributors, supplied problems.json, record 3163 (OPG-139), OpenGarden collection. Catalogue identity and source transcription; CC BY 4.0.
\url{https://huggingface.co/datasets/ulamai/UnsolvedMath}.
\item[S2] Open Problem Garden, (m,n)-cycle covers. Original problem and discussion; the mathematical formulation below is expanded from the supplied source record.
\url{http://www.openproblemgarden.org/op/m_n_cycle_covers}.
\end{itemize}
\end{document}
